\section{Circuit Complexity --- Feed Forward}
%%% generated by the blueprint pipeline from TCSlib.Complexity.CircuitComplexity.LayeredCircuit
%%%%%%%%%%%%%%%%%%%%%%%%%%%%%%%%%%%%%%%%%%%%%%%%%%%%%%%%%%%%%%%%%%%%%%%%%%%%%%%
\subsection*{Circuit Complexity --- Feed Forward: Introduction}

This file introduces layered feedforward circuits over an arbitrary alphabet $\alpha$:
gates with explicit input wiring, the evaluation of nodes and of the whole circuit,
and the basic complexity measures (depth, size, finiteness, gate set restriction).
It then relates feedforward circuits to the tree-shaped \texttt{BoolCircuit.TreeCircuit}
in both directions: tree-unrolling a layered AND/OR circuit into a tree (with an
exponential-in-depth size bound), and embedding a tree circuit back as a
feedforward circuit padded with identity wires.

%%%%%%%%%%%%%%%%%%%%%%%%%%%%%%%%%%%%%%%%%%%%%%%%%%%%%%%%%%%%%%%%%%%%%%%%%%%%%%%
\subsection*{Circuit Complexity --- Feed Forward: Results}

\begin{definition}[Gate with input wiring]
\lean{BoolCircuit.Gate}
\statementsource{lmn-constant-depth-fourier}{pdf-63759bd909d5-p004-b007}
\statementsource{goldreich-approximation-method}{
  pdf-988dc9f0ae91-p004-b004,
  pdf-988dc9f0ae91-p005-b006
}
\label{BoolCircuit.Gate}
\leanok
\uses{BoolCircuit.GateOp}
A gate over alphabet $\alpha$ with input domain \texttt{domain} consists of a gate
operation \texttt{op} (an arity type $\iota$ together with a function
$(\iota \to \alpha) \to \alpha$) and a wiring map $\iota \to \texttt{domain}$ that
assigns to each input slot the node it reads from.
\end{definition}

\begin{definition}[Layered feedforward circuit]
\lean{BoolCircuit.LayeredCircuit}
\statementsource{lmn-constant-depth-fourier}{pdf-63759bd909d5-p004-b007}
\statementsource{goldreich-approximation-method}{
  pdf-988dc9f0ae91-p004-b004,
  pdf-988dc9f0ae91-p005-b002
}
\label{BoolCircuit.LayeredCircuit}
\leanok
\uses{BoolCircuit.Gate}
A feedforward circuit from \texttt{inp} to \texttt{out} over $\alpha$ consists of a
depth $d \in \bbn$, a family of node types indexed by the layers
$\mathrm{Fin}\,(d+1)$, and, for every layer $i < d$, a gate for each node of layer
$i+1$ whose inputs are wired to nodes of layer $i$. Layer $0$ is required to be the
input type \texttt{inp} and the last layer the output type \texttt{out}.
\end{definition}

\begin{definition}[Identity gate operation]
\lean{BoolCircuit.LayeredCircuit.GateOp.id}
\label{BoolCircuit.LayeredCircuit.GateOp.id}
\leanok
\uses{BoolCircuit.GateOp}
The gate operation of arity $\mathrm{PUnit}$ that returns its single input
unchanged.
\end{definition}

\begin{definition}[Evaluation of a single gate]
\lean{BoolCircuit.LayeredCircuit.Gate.eval}
\label{BoolCircuit.LayeredCircuit.Gate.eval}
\leanok
\uses{BoolCircuit.Gate}
Given a gate $g$ and an assignment $xs$ of values to the nodes of its input domain,
the value of $g$ is $g.\texttt{op}.\texttt{func}\,(xs \circ g.\texttt{inputs})$, i.e.\
the gate operation applied to the values read along the input wires.
\end{definition}

\begin{definition}[Evaluation of a node]
\lean{BoolCircuit.LayeredCircuit.evalNode}
\label{BoolCircuit.LayeredCircuit.evalNode}
\leanok
\uses{BoolCircuit.LayeredCircuit.Gate.eval}
For an input assignment $xs : \texttt{inp} \to \alpha$, the value of a node at layer
$d$ is defined by recursion on $d$: a layer-$0$ node is read off directly from $xs$,
and a node at layer $n+1$ evaluates its gate on the values of the layer-$n$ nodes.
\end{definition}

\begin{definition}[Evaluation of a circuit]
\lean{BoolCircuit.LayeredCircuit.eval}
\label{BoolCircuit.LayeredCircuit.eval}
\leanok
\uses{BoolCircuit.LayeredCircuit.evalNode}
The function $\texttt{out} \to \alpha$ sending each output node to its value in the
last layer under the input assignment $xs$.
\end{definition}

\begin{definition}[Evaluation with a unique output]
\lean{BoolCircuit.LayeredCircuit.eval₁}
\label{BoolCircuit.LayeredCircuit.eval₁}
\leanok
\uses{BoolCircuit.LayeredCircuit.eval}
When the output type is a \texttt{Unique} type, the single value
$F.\texttt{eval}\,xs\,\texttt{default} \in \alpha$ computed by the circuit.
\end{definition}

\begin{definition}[Size of a feedforward circuit]
\lean{BoolCircuit.LayeredCircuit.size}
\statementsource{lmn-constant-depth-fourier}{pdf-63759bd909d5-p004-b007}
\statementsource{goldreich-approximation-method}{
  pdf-988dc9f0ae91-p004-b004,
  pdf-988dc9f0ae91-p005-b002
}
\label{BoolCircuit.LayeredCircuit.size}
\leanok
The size of $F$ is the number of its non-input nodes, i.e.\ the total number of
nodes lying in layers $1, \dots, \mathrm{depth}(F)$.
\end{definition}

\begin{definition}[Finiteness of all layers]
\lean{BoolCircuit.LayeredCircuit.Finite}
\label{BoolCircuit.LayeredCircuit.Finite}
\leanok
The proposition that the node type of every layer $i$ is finite.
\end{definition}

\begin{definition}[Restriction to a gate set]
\lean{BoolCircuit.LayeredCircuit.onlyUsesGates}
\statementsource{goldreich-approximation-method}{
  pdf-988dc9f0ae91-p004-b004,
  pdf-988dc9f0ae91-p005-b002,
  pdf-988dc9f0ae91-p005-b006
}
\label{BoolCircuit.LayeredCircuit.onlyUsesGates}
\leanok
\uses{BoolCircuit.GateOp}
For a set $S$ of gate operations, $\texttt{onlyUsesGates}\ S$ asserts that the
operation of every gate of the circuit belongs to $S$.
\end{definition}

\begin{theorem}[Positive circuit size]
\lean{BoolCircuit.TreeCircuit.one_le_size}
\label{BoolCircuit.TreeCircuit.one_le_size}
\leanok
\uses{BoolCircuit.TreeCircuit, BoolCircuit.TreeCircuit.size}
\difficulty{2}
For every natural number $n$ and every Boolean circuit $c$ on $n$ variables, its size is
at least $1$.
\end{theorem}

\begin{definition}[AND/OR gate restriction]
\lean{BoolCircuit.LayeredCircuit.IsAndOrGate}
\statementsource{lmn-constant-depth-fourier}{pdf-63759bd909d5-p004-b007}
\statementsource{goldreich-approximation-method}{
  pdf-988dc9f0ae91-p004-b004,
  pdf-988dc9f0ae91-p005-b006
}
\label{BoolCircuit.LayeredCircuit.IsAndOrGate}
\leanok
\uses{BoolCircuit.LayeredCircuit}
Given a Boolean feedforward circuit $F$, a flag $\texttt{isAnd}$ on each non-input
node and a finiteness witness $\texttt{gfin}$ for each gate's arity,
$F.\texttt{IsAndOrGate}$ asserts that every gate operation is the conjunction of its
inputs when the flag is \texttt{true} and their disjunction when it is \texttt{false},
the inputs being enumerated by $\texttt{gfin}$. This is the gate restriction under
which $F$ can be converted into a \texttt{BoolCircuit.TreeCircuit}.
\end{definition}

\begin{definition}[Tree-unrolling of a node]
\lean{BoolCircuit.nodeToTreeCircuit}
\label{BoolCircuit.nodeToTreeCircuit}
\leanok
\uses{BoolCircuit.LayeredCircuit, BoolCircuit.TreeCircuit}
Recursively expands a node $v$ of layer $m$ of a Boolean feedforward circuit into a
tree-shaped $\texttt{Circuit}\ n$: a layer-$0$ node becomes the positive literal of the
corresponding input variable, and a node at a higher layer becomes an AND/OR node
whose children are the unrollings of the nodes on its input wires. Nodes feeding
several downstream gates are duplicated.
\end{definition}

\begin{theorem}[Evaluation preserved by tree unrolling]
\lean{BoolCircuit.nodeToTreeCircuit_eval}
\label{BoolCircuit.nodeToTreeCircuit_eval}
\leanok
\uses{BoolCircuit.LayeredCircuit, BoolCircuit.LayeredCircuit.Gate.eval, BoolCircuit.LayeredCircuit.IsAndOrGate, BoolCircuit.LayeredCircuit.evalNode, BoolCircuit.Lit.eval, BoolCircuit.TreeCircuit.eval, BoolCircuit.nodeToTreeCircuit}
\difficulty{5}
Let $F$ be a Boolean layered circuit with input type $	ext{Fin }n$, and assign to each
non-input node an AND or OR flag. Suppose every gate has finite arity and computes the
conjunction of its inputs when its flag is AND, and their disjunction otherwise. For
every layer index $m<	ext{depth}(F)+1$, every node $v$ in that layer, and every Boolean
assignment $x$ to the inputs, the tree obtained by recursively unrolling $v$ evaluates
to the same Boolean value as $v$ in $F$ under $x$.
\end{theorem}

\begin{theorem}[Size bound for tree unrolling]
\lean{BoolCircuit.nodeToTreeCircuit_size_le}
\label{BoolCircuit.nodeToTreeCircuit_size_le}
\leanok
\uses{BoolCircuit.LayeredCircuit, BoolCircuit.TreeCircuit.size, BoolCircuit.nodeToTreeCircuit}
\difficulty{6}
Let $F$ be a Boolean layered circuit with input type $	ext{Fin }n$, and assign an
AND-or-OR flag to each non-input node. Suppose every gate has finite arity, and for some
natural number $k$ every gate's arity is at most $k$. For every $m<	ext{depth}(F)+1$ and
every node $v$ at layer $m$, the tree obtained by recursively unrolling $v$ has size at
most $(k+1)^m$.
\end{theorem}

\begin{definition}[Conversion to a Boolean circuit]
\lean{BoolCircuit.LayeredCircuit.toTreeCircuit}
\label{BoolCircuit.LayeredCircuit.toTreeCircuit}
\leanok
\uses{BoolCircuit.LayeredCircuit, BoolCircuit.TreeCircuit, BoolCircuit.nodeToTreeCircuit}
For an AND/OR feedforward circuit $F$ and a chosen output node $o : \texttt{out}$,
the tree-shaped $\texttt{Circuit}\ n$ obtained by tree-unrolling $o$ from the last
layer down to the inputs.
\end{definition}

\begin{theorem}[Tree-unrolling preserves evaluation]
\lean{BoolCircuit.LayeredCircuit.toTreeCircuit_eval}
\label{BoolCircuit.LayeredCircuit.toTreeCircuit_eval}
\leanok
\uses{BoolCircuit.LayeredCircuit, BoolCircuit.LayeredCircuit.IsAndOrGate, BoolCircuit.LayeredCircuit.eval, BoolCircuit.LayeredCircuit.toTreeCircuit, BoolCircuit.TreeCircuit.eval, BoolCircuit.nodeToTreeCircuit_eval}
\difficulty{1}
Let $F$ be a layered Boolean circuit with input nodes indexed by $\mathrm{Fin}\,n$ and
output type $\mathrm{out}$. Suppose each non-input node is designated AND or OR, every
gate has finite arity, and each gate computes the conjunction of its input values when
designated AND and their disjunction when designated OR. Then for every output node $o$
and every Boolean input assignment $x$, the tree circuit obtained by unrolling $o$
evaluates to the same value as $F$ at $o$ under $x$.
\end{theorem}

\begin{theorem}[Size bound for tree-unrolling]
\lean{BoolCircuit.LayeredCircuit.toTreeCircuit_size_le}
\label{BoolCircuit.LayeredCircuit.toTreeCircuit_size_le}
\leanok
\uses{BoolCircuit.LayeredCircuit, BoolCircuit.LayeredCircuit.toTreeCircuit, BoolCircuit.TreeCircuit.size, BoolCircuit.nodeToTreeCircuit_size_le}
\difficulty{1}
Let $F$ be a layered Boolean circuit with input nodes indexed by $\mathrm{Fin}\,n$ and
output type $\mathrm{out}$. Designate each non-input node as AND or OR, and suppose
every gate has finite arity. If there is a natural number $k$ such that every gate has
at most $k$ inputs, then for every output node $o$, the tree circuit obtained by
unrolling $o$ has size at most $(k+1)^{\mathrm{depth}(F)}$.
\end{theorem}

\begin{definition}[Embedding a Boolean circuit as a feedforward circuit]
\lean{BoolCircuit.TreeCircuit.toLayeredWrapper}
\label{BoolCircuit.TreeCircuit.toLayeredWrapper}
\leanok
\uses{BoolCircuit.LayeredCircuit, BoolCircuit.LayeredCircuit.GateOp.id, BoolCircuit.TreeCircuit, BoolCircuit.TreeCircuit.depth, BoolCircuit.TreeCircuit.eval}
Turns a circuit $C$ on $n$ inputs into a feedforward circuit
$\texttt{FeedForward}\ \texttt{Bool}\ (\mathrm{Fin}\,n)\ \mathrm{Unit}$ of depth
$C.\texttt{depth} + 1$: layer $0$ carries the $n$ input variables and every higher
layer carries a single wire. The gate from layer $0$ to layer $1$ computes
$C.\texttt{eval}$ from all inputs at once, and all later gates are identity wires
that pass the single Boolean value upward.
\end{definition}

\begin{theorem}[Evaluation of embedded circuit nodes]
\lean{BoolCircuit.TreeCircuit.toLayeredWrapper_evalNode_const}
\label{BoolCircuit.TreeCircuit.toLayeredWrapper_evalNode_const}
\leanok
\uses{BoolCircuit.LayeredCircuit.evalNode, BoolCircuit.TreeCircuit, BoolCircuit.TreeCircuit.depth, BoolCircuit.TreeCircuit.eval, BoolCircuit.TreeCircuit.toLayeredWrapper}
\difficulty{4}
Let $C$ be a Boolean circuit on $n$ variables and let $x : \mathrm{Fin}\,n \to \bbf_2$
be an input assignment. For every integer $m$ with $1 \le m \le \mathrm{depth}(C)+1$ and
every node $v$ in layer $m$ of the feedforward circuit embedding of $C$, the value of
$v$ on input $x$ equals the value of $C$ on $x$.
\end{theorem}

\begin{theorem}[Evaluation of the layered embedding]
\lean{BoolCircuit.TreeCircuit.toLayeredWrapper_eval}
\label{BoolCircuit.TreeCircuit.toLayeredWrapper_eval}
\leanok
\uses{BoolCircuit.LayeredCircuit.eval₁, BoolCircuit.TreeCircuit, BoolCircuit.TreeCircuit.eval, BoolCircuit.TreeCircuit.toLayeredWrapper, BoolCircuit.TreeCircuit.toLayeredWrapper_evalNode_const}
\difficulty{2}
Let $C$ be a Boolean circuit tree on $n$ variables, and let $x$ assign a Boolean value
to each variable. The value computed by the unique output of its layered feedforward
embedding on input $x$ equals the value obtained by evaluating $C$ directly on $x$.
\end{theorem}

\begin{theorem}[Depth of the layered embedding]
\lean{BoolCircuit.TreeCircuit.toLayeredWrapper_depth}
\label{BoolCircuit.TreeCircuit.toLayeredWrapper_depth}
\leanok
\uses{BoolCircuit.TreeCircuit, BoolCircuit.TreeCircuit.depth, BoolCircuit.TreeCircuit.toLayeredWrapper}
\difficulty{1}
Let $C$ be a Boolean circuit tree on $n$ variables. Its layered feedforward embedding
has depth exactly one more than the depth of $C$.
\end{theorem}

\begin{theorem}[Size bound for the circuit embedding]
\lean{BoolCircuit.TreeCircuit.toLayeredWrapper_size_le}
\label{BoolCircuit.TreeCircuit.toLayeredWrapper_size_le}
\leanok
\uses{BoolCircuit.LayeredCircuit.size, BoolCircuit.TreeCircuit, BoolCircuit.TreeCircuit.depth, BoolCircuit.TreeCircuit.one_le_size, BoolCircuit.TreeCircuit.size, BoolCircuit.TreeCircuit.toLayeredWrapper}
\difficulty{3}
For every Boolean circuit $C$ on $n$ variables, the number of non-input nodes in its
feedforward circuit embedding is at most $\mathrm{size}(C)(\mathrm{depth}(C)+1)$.
\end{theorem}

%%% added by the blueprint pipeline
%%%%%%%%%%%%%%%%%%%%%%%%%%%%%%%%%%%%%%%%%%%%%%%%%%%%%%%%%%%%%%%%%%%%%%%%%%%%%%%
\subsection*{Circuit Complexity --- Feed Forward: Further results}

\begin{definition}[Gate operation]
\lean{BoolCircuit.GateOp}
\statementsource{goldreich-approximation-method}{
  pdf-988dc9f0ae91-p004-b004,
  pdf-988dc9f0ae91-p005-b006
}
\label{BoolCircuit.GateOp}
\leanok
A single operation in a feedforward circuit over an alphabet $\alpha$: it consists
of an arity type $\iota$ together with a function $(\iota \to \alpha) \to \alpha$
that computes the gate's output value from the values supplied on its input slots.
\end{definition}
